\usepackage{newunicodechar}
\newunicodechar{₀}{\textsubscript{0}}
\newunicodechar{₁}{\textsubscript{1}}
\newunicodechar{₂}{\textsubscript{2}}
\newunicodechar{₃}{\textsubscript{3}}
\newunicodechar{ᵢ}{\textsubscript{i}}
\newunicodechar{ⱼ}{\textsubscript{j}}
\newunicodechar{ₖ}{\textsubscript{k}}
\newunicodechar{σ}{\ensuremath{\sigma}}
\newunicodechar{μ}{\ensuremath{\mu}}
\newunicodechar{ε}{\ensuremath{\varepsilon}}
\newunicodechar{α}{\ensuremath{\alpha}}
\newunicodechar{β}{\ensuremath{\beta}}
\newunicodechar{γ}{\ensuremath{\gamma}}
\newunicodechar{λ}{\ensuremath{\lambda}}
% APA 7: texto alineado a la izquierda (flush left, ragged right)
\usepackage{ragged2e}
\raggedright
\setlength{\emergencystretch}{3em}
% APA 7: interlineado doble
\usepackage{setspace}
\doublespacing
% APA 7: sangría de 0.5 pulgadas en primera línea de párrafo
\setlength{\parindent}{0.5in}
\setlength{\parskip}{0pt}
% APA 7: sangría francesa (hangindent) para las entradas de referencias
\newcommand{\hangref}{\par\parindent=0pt\leavevmode\hangindent=0.5in\hangafter=1}
